% Outline only. References/appendices are outside the eight-page main-text budget.
% Source mapping: content_mapping.md, Appendices A--E.

\section{Policy-generation prompts and interfaces}
\label{app:generation}

\subsection{Shared generation specification}
\label{app:specification}
\begin{itemize}
    \item Reuse the IJCAI26 policy-generation specification: semantic variables,
    differentiable computation, parameter initialization, action-distribution
    outputs, and the logging dictionary. Preserve prompts used in experiments;
    distinguish historical prompt text from any corrected explanatory text.
\end{itemize}

\subsection{Environment descriptions and observation schemas}
\label{app:environment}
\begin{itemize}
    \item Reuse the Car Racing and Door Opening environment prompts, feature
    shapes, sign conventions, action spaces, and task information. Verify these
    against the actual experiment interface before drafting.
\end{itemize}

\section{Diagnostic and revision prompts}
\label{app:prompts}

\subsection{Hypothesis generation and analysis execution}
\label{app:diagnostic-prompts}
\begin{itemize}
    \item Include the shared diagnostic system prompt, policy/schema input
    template, result format, and execution details. Show full scripts for the
    main example and document error handling from actual implementation records.
\end{itemize}

\subsection{Revision, additional tests, and termination}
\label{app:revision-prompts}
\begin{itemize}
    \item Reuse all three response branches from the refinement prompt. Add
    the actual orchestration, budget, stopping, and any selection rules from
    experiment records; identify behavior not documented by the old paper.
\end{itemize}

\section{Training and evaluation details}
\label{app:training}

\subsection{Imitation learning and refinement configuration}
\label{app:il}
\begin{itemize}
    \item Provide optimizer settings, initialization/transfer behavior, training
    stopping criteria, LLM configuration, rollout seeds, evaluation separation,
    independent-run counts, and uncertainty calculations from saved records.
\end{itemize}

\subsection{Post-refinement reinforcement learning}
\label{app:rl}
\begin{itemize}
    \item Adapt the IJCAI26 RL subsection after checking objective, sequence
    grouping, indexing, and implementation consistency. State the assumption
    behind grouping states at corresponding times across sampled trajectories.
    Document RL hyperparameters and exact reset accounting.
\end{itemize}

\section{Complete refinement records}
\label{app:records}

\subsection{Car Racing programs, diagnostics, and rollouts}
\label{app:car-record}
\begin{itemize}
    \item Recover the before/after programs, demonstration-fitting details,
    relevant log rows, full diagnostic outputs, and post-refit evaluations for
    the main example. Record provenance and iteration identifiers.
\end{itemize}

\subsection{Door Opening progression}
\label{app:door-record}
\begin{itemize}
    \item Expand the existing iteration images and narrative with matched
    programs/diagnostics where available. Preserve ineffective intermediate
    revisions and distinguish measured distance from qualitative task success.
\end{itemize}

\section{Extended feedback and semantic analyses}
\label{app:extended}

\subsection{Diagnostic and strategy coverage}
\label{app:coverage}
\begin{itemize}
    \item Reuse the detailed diagnostic-coverage and semantic-strategy tables.
    Explain manual categories, coding criteria, and limits of coverage counts
    as measures of semantic correctness or diagnostic usefulness.
\end{itemize}

\subsection{Feedback budgets and supplementary results}
\label{app:budgets}
\begin{itemize}
    \item Document raw-table subsampling, feedback comparability, token
    aggregation, refinement histories, and supplementary baseline/RL rows from
    existing records. Do not populate missing measurements with inferred data.
\end{itemize}
